\documentclass[11pt,a4paper]{article}
\usepackage{fontspec,amsmath,amssymb,amsthm,microtype}
\setmainfont{lmroman10-regular.otf}[BoldFont=lmroman10-bold.otf,ItalicFont=lmroman10-italic.otf,BoldItalicFont=lmroman10-bolditalic.otf]
\usepackage{xcolor}
\usepackage[most]{tcolorbox}
\usepackage{mathrsfs}
\providecommand{\lt}{<}
\providecommand{\gt}{>}
\setmonofont{DejaVuSansMono.ttf}[BoldFont=DejaVuSansMono-Bold.ttf,ItalicFont=DejaVuSansMono-Oblique.ttf,BoldItalicFont=DejaVuSansMono-BoldOblique.ttf,Scale=0.86]
\newfontfamily\leanfallback{STIXTwoMath-Regular.otf}[Scale=0.86]
\definecolor{leanrule}{gray}{0.62}
\definecolor{leanback}{gray}{0.962}
\definecolor{leantitle}{gray}{0.22}
\definecolor{leanmuted}{gray}{0.38}
\newtcolorbox{leanbox}[1]{enhanced,breakable,colback=leanback,colframe=leanrule,
  boxrule=0.4pt,arc=0.6mm,left=2.4mm,right=2.4mm,top=1.8mm,bottom=1.8mm,
  before skip=0.8em,after skip=1em,fonttitle=\small\bfseries,coltitle=leantitle,
  colbacktitle=leanback,titlerule=0.3pt,title={#1}}
\newlength{\leancharwidth}
\newenvironment{leancode}{\par\smallskip\begingroup\ttfamily\footnotesize\raggedright
  \settowidth{\leancharwidth}{M}\setlength{\parskip}{0pt}\setlength{\parindent}{0pt}
  \frenchspacing}{\par\endgroup\smallskip}
\newcommand{\leanline}[2]{\hangindent=\dimexpr#1\leancharwidth+4\leancharwidth\relax
  \hangafter=1\noindent\hspace*{#1\leancharwidth}#2\par}
\newcommand{\leansource}[1]{\par{\footnotesize\color{leanmuted}#1}\par}
\newcommand{\leanlabel}[1]{\par\smallskip\noindent{\small\bfseries #1}\par}
\usepackage[margin=25mm]{geometry}
\usepackage[unicode,hidelinks]{hyperref}
\hypersetup{pdfauthor={LangevinLean Project}}
\setlength{\emergencystretch}{3em}
\setlength{\parskip}{0.4em}
\setlength{\parindent}{0pt}
\allowdisplaybreaks[2]
\raggedbottom
\title{\texorpdfstring{From the constructed semigroup corrector to the comparison process}{From the constructed semigroup corrector to the comparison process}}
\author{LangevinLean Project}
\date{}
\begin{document}\maketitle
Draft mathematical exposition; these notes may contain mistakes.


The bounded corrector of a reversible semigroup, its derived real-range preservation and actual stationary Markov evaluations supply a real Poisson comparison process. Corrector norms, Dirichlet domination and the exact integrated source are consequences of the construction.


\section{\texorpdfstring{The genuine source integral and evaluated process}{The genuine source integral and evaluated process}}


\textbf{Construction.} Extract the real corrector \(G_h^{\mathbb R}\) and real transition \(P_h^{\mathbb R}\), and define the source operator by \((I-P_h^{\mathbb R})G_h^{\mathbb R}\). The already proved Poisson identity identifies this operator with \(f\mapsto\int_0^h(f-P_t^{\mathbb R}f)\,dt\).



\textbf{Proof.} The real-extraction theorems provide the exact integral identity, the norm bound \(h\) and energy factor \(h^2\). Combine them with the supplied actual evaluation isometries, their adaptation and their conditional Markov identity. The resulting process uses the constructed corrector itself and retains all these proved estimates.

\begin{leanbox}{Lean: \texttt{semigroupPoissonSource}, \texttt{semigroupPoissonSource\_\allowbreak{}eq\_\allowbreak{}integral}, \texttt{semigroupPoissonProcess}}
\leansource{\texttt{semigroupPoissonSource} \quad WassersteinMixing/\allowbreak{}SemigroupProcess.lean, lines 26–28}
\begin{leancode}
\leanline{0}{def\ semigroupPoissonSource\ (S\ :\ ReversibleSemigroup\ H)\ (i\ :\ R\ →ₗᵢ[ℝ]\ H)}
\leanline{4}{(h\ :\ ℝ)\ (hh\ :\ 0\ ≤\ h)\ :\ R\ →L[ℝ]\ R\ :=}
\leanline{2}{realCorrector\ S\ i\ h\ hh\ -\ realOp\ S\ i\ h\ ∘L\ realCorrector\ S\ i\ h\ hh}
\end{leancode}
\leansource{\texttt{semigroupPoissonSource\_\allowbreak{}eq\_\allowbreak{}integral} \quad WassersteinMixing/\allowbreak{}SemigroupProcess.lean, lines 30–34}
\begin{leancode}
\leanline{0}{theorem\ semigroupPoissonSource\_eq\_integral\ (S\ :\ ReversibleSemigroup\ H)}
\leanline{4}{(i\ :\ R\ →ₗᵢ[ℝ]\ H)\ (hreal\ :\ PreservesReal\ S\ i)\ (h\ :\ ℝ)\ (hh\ :\ 0\ ≤\ h)\ (f\ :\ R)\ :}
\leanline{4}{semigroupPoissonSource\ S\ i\ h\ hh\ f\ =\ ∫\ t\ in\ 0..h,\ f\ -\ realOp\ S\ i\ t\ f}
\end{leancode}
\leansource{\texttt{semigroupPoissonProcess} \quad WassersteinMixing/\allowbreak{}SemigroupProcess.lean, lines 38–54}
\begin{leancode}
\leanline{0}{def\ semigroupPoissonProcess\ (Q\ :\ Measure\ Ω)\ (S\ :\ ReversibleSemigroup\ H)}
\leanline{4}{(i\ :\ R\ →ₗᵢ[ℝ]\ H)\ (hreal\ :\ PreservesReal\ S\ i)\ (h\ :\ ℝ)\ (hh\ :\ 0\ <\ h)}
\leanline{4}{(filtration\ :\ Filtration\ ℕ\ mΩ)\ (eval\ :\ ℕ\ →\ R\ →ₗᵢ[ℝ]\ Lp\ E\ 2\ Q)}
\leanline{4}{(hadapted\ :\ ∀\ n\ f,\ AEStronglyMeasurable[filtration\ n]\ (eval\ n\ f)\ Q)}
\leanline{4}{(hmarkov\ :\ ∀\ n\ f,}
\leanline{6}{(condExpL2\ E\ ℝ\ (filtration.le\ n)\ (eval\ (n\ +\ 1)\ f)\ :\ Lp\ E\ 2\ Q)\ =}
\leanline{8}{eval\ n\ (realOp\ S\ i\ h\ f))\ :\ PoissonProcess\ R\ E\ Q\ h\ hh\ where}
\leanline{2}{transition\ :=\ realOp\ S\ i\ h}
\leanline{2}{corrector\ :=\ realCorrector\ S\ i\ h\ hh.le}
\leanline{2}{poissonSource\ :=\ semigroupPoissonSource\ S\ i\ h\ hh.le}
\leanline{2}{corrector\_norm\ :=\ realCorrector\_norm\_le\ S\ i\ hreal\ h\ hh.le}
\leanline{2}{corrector\_energy\ :=\ realCorrector\_energy\_le\ S\ i\ hreal\ h\ hh.le\ hh.le}
\leanline{2}{poisson\_identity\ :=\ fun\ \_\ ↦\ rfl}
\leanline{2}{filtration\ :=\ filtration}
\leanline{2}{eval\ :=\ eval}
\leanline{2}{adapted\ :=\ hadapted}
\leanline{2}{markov\ :=\ hmarkov}
\end{leancode}
\leanlabel{How the Lean proof follows the argument}
\texttt{semigroup\allowbreak{}Poisson\allowbreak{}Source} is the bounded operator difference. \texttt{semigroup\allowbreak{}Poisson\allowbreak{}Source\_\allowbreak{}eq\_\allowbreak{}integral} invokes \texttt{real\allowbreak{}Corrector\_\allowbreak{}poisson}. \texttt{semigroup\allowbreak{}Poisson\allowbreak{}Process} fills every corrector field from the proved real-extraction theorems, leaving only actual evaluation and Markov facts to the concrete process construction.
\end{leanbox}

\end{document}
